
\documentclass{article}
\usepackage{colortbl}
\usepackage{makecell}
\usepackage{multirow}
\usepackage{supertabular}

\begin{document}

\newcounter{utterance}

\centering \large Interaction Transcript for game `eqbench', experiment `eq\_bench', episode 52 with reflexion{-}Qwen3.5{-}27B.
\vspace{24pt}

{ \footnotesize  \setcounter{utterance}{1}
\setlength{\tabcolsep}{0pt}
\begin{supertabular}{c@{$\;$}|p{.15\linewidth}@{}p{.15\linewidth}p{.15\linewidth}p{.15\linewidth}p{.15\linewidth}p{.15\linewidth}}
    \# & \multicolumn{2}{c}{Player} && \multicolumn{2}{c}{Game Master} \\
    \hline

    \theutterance \stepcounter{utterance}  
    & & & \multicolumn{4}{p{0.6\linewidth}}{
        \cellcolor[rgb]{0.9,0.9,0.9}{
            \makecell[{{p{\linewidth}}}]{
                \texttt{\tiny{[P1$\langle$GM]}}
                \texttt{You have attempted this same task before and did not succeed. These are the plans you wrote after those attempts.} \\
\texttt{Trial \#1: New plan: My previous attempt failed because I assigned a high score to "Challenged" (9) and a moderate score to "Angry" (6), likely misinterpreting the detective's tone. Ben's dialogue ("inevitable," "not as invincible," "can't keep this up") suggests a tone of absolute certainty and professional resolve rather than being personally challenged or merely angry. The most accurate emotional response for a detective who believes they have the upper hand and sees the criminal's defiance as futile is "Proud" (in their professional capability) or a very specific, cold "Angry" at the criminal's arrogance, but definitely not "Challenged" as he feels he has already won. Next time, I will re{-}evaluate the power dynamic: Ben is not being challenged; he is asserting dominance. I will lower the "Challenged" score to near zero and increase the "Angry" score to reflect his frustration with Ace's lack of remorse, or potentially "Proud" if the context implies satisfaction in the inevitable capture. Specifically, I will prioritize "Angry" as the primary emotion due to the moral weight of the theft ("people you're stealing from") and the criminal's dismissive attitude, while ensuring "Challenged" is scored low (0{-}2) because Ben explicitly states the outcome is "inevitable," implying he feels no doubt or challenge.} \\
\\ 
\texttt{Trial \#2: New plan: My previous attempt failed because I assigned a high score to "Angry" (8) while underestimating "Challenged" (1), misinterpreting the detective's reaction to Ace's direct insult. While Ben is certainly frustrated by the theft, the dialogue culminates in Ace explicitly attacking Ben's competence ("you're not as clever as you think you are"). A detective's primary reaction to a direct challenge to their intelligence and capability is to feel "Challenged," not just angry. Ben's earlier statements about the "inevitable" capture were a setup for this final confrontation; Ace's retort shifts the dynamic from a one{-}sided pursuit to a personal test of wits. Next time, I will prioritize "Challenged" as the highest score (8{-}9) because the dialogue ends on a direct provocation of Ben's ego and skill. I will lower "Angry" to a moderate score (4{-}5) to reflect his professional frustration rather than uncontrolled rage, and keep "Humiliated" low (0{-}1) as Ben maintains his composure. "Proud" should remain low (0{-}2) as the situation is tense, not a moment of triumph. Specifically, I will ensure the score for "Challenged" is the highest to reflect the final line's impact.} \\
\\ 
\texttt{Trial \#3: New plan: My previous attempts failed because I oscillated between prioritizing "Challenged" and "Angry" without fully accounting for the cumulative effect of the dialogue's moral stakes. In Trial 1, I underestimated the impact of Ace's final insult on Ben's ego. In Trial 2, I overcorrected by making "Challenged" the sole focus, ignoring the deep{-}seated frustration Ben expresses regarding the victims ("people you're stealing from") and Ace's dismissive attitude ("They always do"). The correct strategy is to recognize that Ben feels a complex mix: he is **Angry** (high score) due to the moral violation and Ace's arrogance, and **Challenged** (high score) due to the direct attack on his competence, but these are distinct drivers. "Proud" is unlikely as the situation is unresolved and tense, and "Humiliated" is too strong as Ben maintains a professional, albeit frustrated, demeanor. Next time, I will assign high scores to both "Angry" and "Challenged" (e.g., 8 and 9 respectively) to reflect the dual nature of the conflict: the moral outrage of the theft and the personal provocation of the final line. I will ensure "Proud" and "Humiliated" remain low (0{-}2) to accurately reflect the lack of triumph or total loss of dignity. Specifically, I will score "Challenged" as the highest (9) because the dialogue ends on a direct personal attack, but I will keep "Angry" very high (8) to acknowledge the sustained moral conflict, rather than treating them as mutually exclusive.} \\
\\ 
\\ 
\\ 
\texttt{Your task is to predict the likely emotional responses of a character in this dialogue:} \\
\\ 
\texttt{Ben: You're playing a dangerous game, Ace.} \\
\texttt{Ace: Isn't that the point, Detective? Life's no fun without a little risk.} \\
\texttt{Ben: This isn't a game to the people you're stealing from.} \\
\texttt{Ace: Oh, they'll survive. They always do.} \\
\texttt{Ben: And what about you, Ace? Will you survive when I catch up to you?} \\
\texttt{Ace: You seem pretty confident about that.} \\
\texttt{Ben: I'm just stating the inevitable.} \\
\texttt{Ace: Is that so? Or is it just wishful thinking?} \\
\texttt{Ben: You can't keep this up forever, Ace.} \\
\texttt{Ace: Maybe not. But I'll make sure it's a hell of a ride.} \\
\texttt{Ben: You're not as invincible as you think you are.} \\
\texttt{Ace: And you're not as clever as you think you are, Detective.} \\
\texttt{{[}End dialogue{]}} \\
\\ 
\texttt{At the end of this dialogue, Ben would feel...} \\
\texttt{Humiliated} \\
\texttt{Angry} \\
\texttt{Proud} \\
\texttt{Challenged} \\
\\ 
\texttt{Give each of these possible emotions a score from 0{-}10 for the relative intensity that they are likely to be feeling each.} \\
\\ 
\texttt{You must output in the following format, including headings (of course, you should give your own scores), with no additional commentary:} \\
\\ 
\texttt{Humiliated: <score>} \\
\texttt{Angry: <score>} \\
\texttt{Proud: <score>} \\
\texttt{Challenged: <score>} \\
\\ 
\\ 
\texttt{{[}End of answer{]}} \\
\\ 
\texttt{Remember: zero is a valid score, meaning they are likely not feeling that emotion. You must score at least one emotion > 0.} \\
\\ 
\texttt{Your answer:} \\
            }
        }
    }
     \\ \\

    \theutterance \stepcounter{utterance}  
    & \multicolumn{4}{p{0.6\linewidth}}{
        \cellcolor[rgb]{0.561,0.737,0.561}{
            \makecell[{{p{\linewidth}}}]{
                \texttt{\tiny{[P1$\rangle$GM]}}
                \texttt{Humiliated: 1} \\
\texttt{Angry: 8} \\
\texttt{Proud: 2} \\
\texttt{Challenged: 9} \\
            }
        }
    }
    & & \\ \\

    \theutterance \stepcounter{utterance}  
    & & & \multicolumn{2}{p{0.3\linewidth}}{
        \cellcolor[rgb]{0.9,0.9,0.9}{
            \makecell[{{p{\linewidth}}}]{
                \texttt{\tiny{[GM$|$GM]}}
                \texttt{Humiliated: 1} \\
\texttt{Angry: 8} \\
\texttt{Proud: 2} \\
\texttt{Challenged: 9} \\
            }
        }
    }
    & & \\ \\

    \theutterance \stepcounter{utterance}  
    & & & \multicolumn{2}{p{0.3\linewidth}}{
        \cellcolor[rgb]{0.9,0.9,0.9}{
            \makecell[{{p{\linewidth}}}]{
                \texttt{\tiny{[GM$|$GM]}}
                \texttt{Humiliated: 2} \\
\texttt{Angry: 6} \\
\texttt{Proud: 1} \\
\texttt{Challenged: 8} \\
            }
        }
    }
    & & \\ \\

    \theutterance \stepcounter{utterance}  
    & & & \multicolumn{2}{p{0.3\linewidth}}{
        \cellcolor[rgb]{0.9,0.9,0.9}{
            \makecell[{{p{\linewidth}}}]{
                \texttt{\tiny{[GM$|$GM]}}
                \texttt{game\_result = LOSE} \\
            }
        }
    }
    & & \\ \\

    \theutterance \stepcounter{utterance}  
    & & \multicolumn{4}{p{0.6\linewidth}}{
        \cellcolor[rgb]{0.95,0.95,0.95}{
            \makecell[{{p{\linewidth}}}]{
                \texttt{\tiny{[P1$\leftrightarrow$P]}}
                \texttt{Post{-}episode Reflexion (not sent to GM):} \\
\texttt{New plan: My previous attempts failed because I treated "Angry" and "Challenged" as competing primary emotions, leading to inconsistent scoring strategies across trials. I failed to recognize that Ben's dialogue is defined by a specific progression: he starts with moral outrage (Angry) regarding the victims, but the conversation shifts entirely to a personal contest of wits by the end. The final line ("you're not as clever as you think you are") is a direct attack on Ben's professional identity, which overrides the earlier moral anger in terms of immediate emotional impact. A detective who believes the outcome is "inevitable" (as Ben stated) is not "Challenged" in the sense of doubt, but rather "Challenged" in the sense of being provoked to prove his superiority. However, the most dominant emotion for a detective facing a criminal who dismisses the victims and mocks his intelligence is **Angry** (frustration at the lack of remorse and the personal insult), while **Challenged** is a secondary, reactive state. I will correct this by assigning the highest score to **Angry** (9) to reflect the cumulative weight of the moral violation and the personal insult, and a high but secondary score to **Challenged** (7) to acknowledge the provocation. I will ensure **Proud** and **Humiliated** remain low (0{-}1) as Ben is neither triumphant nor defeated. Specifically, I will prioritize the moral dimension of the dialogue ("people you're stealing from") as the root cause of Ben's emotional state, making "Angry" the definitive primary score.} \\
            }
        }
    }
    & \\ \\

\end{supertabular}
}

\end{document}
